%-\
%-> 中文摘要与英文摘要
%->  ucasthesis 类的摘要不用 abstract 环境，而是用 \chapter* + \intobmk
%->  生成书签，并配合 \keywords / \KEYWORDS 输出关键词行。
%-\
\intobmk\chapter*{摘\quad 要}
\setcounter{page}{1}
\pagenumbering{Roman}

虚拟化试图让客户机获得整台机器的使用接口，却始终把对真实资源的最终控制权保留给虚拟化管理层。本文围绕这一核心问题，系统综述从软件模拟到硬件辅助再到嵌套虚拟化的实现路径，范围限定于系统级整机虚拟化。全文以CPU、内存、中断、设备四维为统一框架，逐维回答资源与状态由谁控制、关键路径如何执行、何时介入以及代价何在。软件一侧以QEMU的TCG为例，沿RISC-V到TCG IR再到x86-64的翻译管线，说明guest指令如何在没有硬件支持时被翻译执行，以及翻译块缓存、寄存器分配与helper调用各自的代价；硬件一侧以Armv8.6-A与Linux KVM为例，展示异常级、两阶段地址翻译、GIC与SMMU如何把高频路径的软件介入下放给处理器与系统结构，管理策略与设备语义仍留给软件。嵌套虚拟化分两部分讨论：Turtles通过合成VMCS与多维分页把多层语义组合到单层硬件支持上，代价是每次跨层转发都要多付出世界切换；HyperTurtle以受控的eBPF hyperupcall把L1定义的部分处理逻辑下放到L0就地执行、并通过共享状态协议和回退路径协调两层管理。两部分构成从“正确实现”到“削减往返”的先后关系，而不是两种可以横向比较性能高低的方案。综述判断，虚拟化演进并非软件被硬件简单替换，而是不断调整执行、状态与策略的放置位置；能否在减少跨层往返的同时维持资源控制、状态一致性与上层管理权，才是方案之间真正可比的问题。

\keywords{虚拟化，虚拟机监视器，硬件辅助虚拟化，嵌套虚拟化，QEMU，KVM}

%-\
%-> 英文摘要
%-\
\intobmk\chapter*{Abstract}

Virtualization aims to give a guest machine the full interface of a computer while never handing it final control over the real resources. This review examines how that goal is met along a path from pure software emulation, through hardware assistance, to nested composition, restricted to system-level full-machine virtualization. A single four-dimensional framework --- CPU, memory, interrupt, and device --- runs through the text, and every dimension answers the same questions: who owns the resource and its state, how the critical path executes, when the hypervisor must intervene, and what that intervention costs. On the software side, QEMU's TCG is followed along its RISC-V to TCG IR to x86-64 pipeline, showing how guest instructions execute without hardware support and what translation-block caching, register allocation, and helper calls cost. On the hardware side, Armv8.6-A and Linux KVM show how exception levels, two-stage address translation, the GIC, and the SMMU move software intervention off the hot path while leaving management policy and device semantics to software. Nested virtualization is discussed in two parts: Turtles composes multi-level semantics onto single-level hardware support by synthesizing VMCS structures and collapsing page tables, paying with additional world switches on every cross-level forward; HyperTurtle uses verified eBPF hyperupcalls to run selected L1 logic in place at L0, coordinating the two levels through a shared-state protocol and a fallback path. The two form a sequence from correctness to cost reduction rather than two alternatives to be ranked by performance. The review concludes that virtualization has not simply replaced software with hardware; it has repeatedly relocated execution, state, and policy. Whether a design can cut cross-level round trips while preserving resource control, state consistency, and upper-level management authority is the question that actually distinguishes one scheme from another.

\KEYWORDS{Virtualization, Hypervisor, Hardware-assisted virtualization, Nested virtualization, QEMU, KVM}

\pagestyle{enfrontmatterstyle}
\cleardoublepage\pagestyle{frontmatterstyle}
